% Part: lambda-calculus
% Chapter: lambda-definability
% Section: introduction

\documentclass[../../../include/open-logic-section]{subfiles}

\begin{document}

\olfileid{lam}{ldf}{int}
\olsection{سريزه}

په لومړي نظر لامبډا حساب د هغو عبارتونو ډېر تجريدي حساب دے چې
تابعې او پر نورو شيانو د هغو تطبيق ښيي۔ د لامبډا حساب په نحو کښې
هېڅ داسې نښه نشته چې وښيي دا د کوم ځانګړي ډول شيانو تابعې دي؛ په
ځانګړې توګه د طبيعي عددونو يادونه پکې نشته۔ د دې حساب بنسټيز عملونه،
يعنې تطبيق او لامبډا تجريد، پر هرې تابع کار کوي، يوازې د طبيعي
عددونو پر تابعو نۀ۔

له دې سره سره، په لږ ځيرکتيا سره د طبيعي عددونو تابعې، يعنې حسابي
تابعې، په لامبډا حساب کښې تعريفولے شو۔ د دې دپاره د هر طبيعي
عدد~$n \in \Nat$ په مقابل کښې يو ځانګړے $\lambd$-ترم~$\num n$
تعريفوو؛ دا د~$n$ \emph{د چرچ عددنښه} ده۔ (د چرچ عددنښې د الونزو
چرچ په نوم نومول شوې دي۔)

\begin{defn}
  که $n \in \Nat$ وي، اړونده \emph{د چرچ عددنښه} $\num{n}$
  د~$n$ استازيتوب کوي:
  \[
    \num{n} \ident \lambd[fx][f^n(x)]
  \]
  دلته $f^n(x)$ د~$x$ پر سر د~$f$ د~$n$ ځلو تطبيق پايله
  ښيي۔ د بېلګې په توګه، $\num{0}$، $\lambd[fx][x]$ او
  $\num{3}$، $\lambd[fx][f(f(f\,x))]$ دے۔
\end{defn}
  
د چرچ عددنښه $\num n$ د داسې لامبډا ترم په توګه کوډ شوې ده چې
دوه آرګومېنټونه $f$ او $x$ اخلي او $f^n(x)$ ورکوي۔ څرګنده ده چې
د چرچ عددنښې په عادي بڼه کښې دي۔

په لامبډا حساب کښې د طبيعي عددونو تمثيل هله ګټور دے چې پرې
محاسبه هم وکړے شو۔ د چرچ له عددنښو سره محاسبه دا معنا لري چې
يو $\lambd$-ترم~$F$ پر داسې عددنښه تطبيق کړو او ګډ ترم~$F\,
\num n$ عادي بڼې ته راکم کړو۔ که هغه تل عادي بڼې ته راکمېږي
او دا بڼه تل د چرچ عددنښه~$\num m$ وي، د محاسبې وتون د~$m$
عدد ګڼلے شو۔ بيا $F$ داسې تابع $f\colon \Nat \to \Nat$ تعريفوي
چې $f(n) = m$ هغه وخت او يوازې هغه وخت وي چې $F\, \num n
\red \num m$ وي۔ د چرچ--روسر د خاصيت له مخې، که عادي بڼه
موجوده وي، يکتا ده۔ نو که $F\, \num n \red \num m$ وي،
$F \, \num n$ بلې عادي بڼې، په ځانګړې توګه بلې د چرچ عددنښې ته،
نۀ شي راکمېدلے۔

پرعکس، که يوه تابع $f\colon \Nat \to \Nat$ راکړل شي، پوښتلے
شو چې آيا داسې ترم $F$ شته چې په دې ډول~$f$ تعريف کړي۔ په دې
حالت کښې وايو $F$،~$f$ \emph{!!{lambda define}s} او $f$،
!!{lambda definable} ده۔ دا خبره څوځاييزو او جزوي تابعو ته هم
غځولے شو۔

\begin{defn}
فرض کړئ $f\colon \Nat^k \to \Nat$۔ وايو يو لامبډا ترم~$F$
هغه وخت $f$ \emph{!!{lambda define}s} چې د ټولو
$n_0$، \dots،~$n_{k-1}$ دپاره،
\[
F \, \num{n_0} \, \num{n_1} \dots \num{n_{k-1}} \red \num{f(n_0, n_1, \dots,
  n_{k-1})}
\]
وي، که $f(n_0, \dots, n_{k-1})$ تعريف شوې وي؛ او که نۀ وي،
$F \, \num{n_0} \,
\num{n_1} \dots \num{n_{k-1}}$ هېڅ عادي بڼه نۀ لري۔
\end{defn}

ثابتې تابعې ډېرې ساده بېلګې دي۔ ترم $C_k \ident
\lambd[x][\num{k}]$ هغه تابع $c_k\colon \Nat \to
\Nat$ !!{lambda define}s چې $c(n) = k$ وي۔ ځکه د هر~$n$
دپاره $C_k \, \num n \ident
(\lambd[x][\num{k}])\num n \redone \num{k}$۔
\textbf{د سرچينې سپيناوی (OLLAM-047):} د ثابتې تابعې نوم
په همدغه جمله کښې له لاندې ارزښت-فورمول سره په يوه نښه
نۀ دے ليکل شوے؛ مطلب هماغه د ټاکلي ارزښت ثابته تابع ده۔
د يووالي تابع $\lambd[x][x]$ !!{lambda defined} کوي۔ د پېچلو تابعو
تعريفول، البته، ډېر ګران دي او ډېر ځله ځيرکتيا غواړي۔ نو دا
ښايي حيرانوونکې وي چې هره محاسبه کېدونکې تابع
!!{lambda definable} ده۔ برعکس خبره هم رښتيا ده: که يوه تابع
!!{lambda definable} وي، نو محاسبه کېدونکې ده۔

\end{document}
